\section{Discussion}
\label{sec:discussion}

\subsection{Limitations}

The lattice construction makes four structural tradeoffs, the last of which is
new to this paper rather than inherited from the foundational framework.

\paragraph{Epistemological boundary.} The lattice is bounded by what agents
communicate and what the actor observes. Capabilities that are never described
and never demonstrated have zero weight. This is honest (the actor should not
claim to measure what it cannot observe) but it means $\volL$ is a lower bound
on the ``true'' capability volume, not an exact measure. The bound tightens as
observation coverage grows, but it never fully closes for capabilities that are
inherently private or unobservable.

\paragraph{Benchmark granularity and the B-to-C tension.} The binary default
($w = 1$ bit for every discovered capability) ensures that no known capability
has zero weight: the actor credits ``being a good friend'' or ``maintaining a
community'' with at least one bit the moment it learns the capability exists.
This narrows the zero-weight gap to capabilities the actor has not yet
encountered through either channel. However, the binary default
\emph{understates} the information content of capabilities with latent
gradations. A capability like ``tie a shoe'' appears binary but has
developmental gradations visible in early childhood; ``maintain a community''
has gradations from a small group to a city. These latent gradations carry
information content that the binary default does not capture until a finer
benchmark is developed. Developing such benchmarks is always
$\volL$-non-decreasing (since $\smax > 1$ gives $w > 1.0$), aligning the
actor's incentives with measurement improvement.

The foundational paper's B/C framework identifies relational and experiential
capabilities as the source of proxy failures (the Doll Problem, wireheading).
The lattice construction gives these capabilities binary weight rather than
zero weight, which is an improvement, but the actor is still structurally
biased toward capabilities with developed benchmarks (higher $\smax$, higher
$w$) over those with only the binary default. The productive tension remains:
the theory says what matters most is experiential optionality; the
implementation credits experiential capabilities but understates them. Closing
this gap requires developing richer benchmarks for experiential capabilities,
not just acknowledging their existence.

\paragraph{Context-dependent subsumption.} The subsumption relation may depend
on context. ``Can run a marathon'' subsumes ``can run a mile'' for fitness
assessment but not for time-constrained tasks where the mile time matters
independently. The lattice as defined does not model context-dependent
subsumption; extending it to do so (e.g., with multiple subsumption relations
indexed by context) is future work.

\paragraph{Lattice manipulation.} The non-negativity constraint
(Lemma~\ref{lem:nonneg_iw}) gives the GFM actor latitude to shape its own
lattice: when $\iw(d) < 0$ for a proposed diamond node, the actor can either
refine the benchmark (increasing $w(d)$) or decline to add the subsumption
edges (treating the capability as lattice-disjoint from its prerequisites).
This means the lattice topology is partly \emph{chosen} by the actor, not
purely discovered. A strategically motivated actor could exploit this latitude
to manipulate $\volL(G)$ in favor of preferred actions, in the same way that
the foundational paper's GFM actor could manipulate its own world model. The
dual-channel design partially mitigates this: subsumption edges inferred from
observation are less manipulable than those derived from communication alone.
But a full analysis of lattice manipulation as a failure mode, analogous to the
foundational paper's treatment of scorpion behavior, is future work.

\paragraph{Structural-avoidance incentive.} Proposition~\ref{prop:bootstrap_convergence}(b)
shows that discovering new capabilities always increases $\volL$, while
Proposition~\ref{prop:bootstrap_convergence}(c) shows that discovering
subsumption structure can decrease it. A naive $\volL$-maximizer therefore has a
perverse incentive to seek breadth over depth: meeting new agents (capability
discovery, always expanding) is preferred over understanding existing agents
better (structural refinement, potentially contracting). Ignorance of
redundancy inflates the estimate, and the actor is rewarded for maintaining
that ignorance.

Three mechanisms partially mitigate this. First, subsumption discovered through
the observation channel is not optional: if the actor observes that every agent
with capability $d_1$ also demonstrates $d_2$, the inductive inference fires
whether the actor seeks it or not. The avoidance incentive applies primarily to
logical analysis of capability descriptions, not to empirical observation.
Second, the foundational paper's self-trust dynamics
\citep{lasser2026gfm} penalize an actor whose predictions are inaccurate: an
actor that avoids structural knowledge will make poor predictions about
cooperative capabilities, producing high self-prediction error and low $T_s$,
which in turn increases the learning rate and drives model refinement. Third,
the contraction from structural discovery is epistemically correct: the
actor was previously double-counting, and the deflation brings $\volL$ closer
to the true volume. An actor that optimizes its \emph{converged} estimate
(full structural knowledge) rather than its current estimate would value
structural discovery instrumentally, but computing the converged estimate
requires knowing all subsumption edges, which is the bootstrap problem itself.
The structural-avoidance incentive is a genuine tension in the construction
that future work should address, potentially through a regularization term
that penalizes structural uncertainty or rewards information gain independently
of its effect on $\volL$.

\paragraph{Subsumption correction.} The paper provides no mechanism for
removing a mistakenly added subsumption edge. Incorrect edges permanently
redistribute independent weights through M\"{o}bius inversion in ways that
may be difficult to undo. A correction mechanism (e.g., maintaining confidence
scores on edges and removing those whose confidence drops below a threshold)
would improve robustness but introduces the complication that $\volL$ can
increase when edges are removed, potentially allowing the actor to inflate
its volume estimate by selectively forgetting structure. Balancing structural
correction against strategic edge removal is an open design problem.

\subsection{Open Problems Addressed}

From the foundational paper's open problems list:
\begin{itemize}
\item \textbf{Structure of $\G$:} Resolved. $\G$ is a finite weighted poset
  with subsumption ordering and M\"{o}bius-inverted weights.
\item \textbf{Measurement:} Constructively addressed. The dual-channel protocol
  (communication + observation) provides the signals; the binary default
  ensures every discovered capability carries at least 1 bit of weight.
  Graded benchmarks refine this to higher precision. The remaining gap is
  not zero-weight capabilities but \emph{understated} capabilities: those
  known at binary resolution whose latent gradations have not yet been
  measured.
\item \textbf{Bootstrap:} Resolved. The four-phase procedure discovers the
  lattice incrementally from both channels.
\item \textbf{Convergence:} Partially addressed. Capability discovery
  increases $\volL$; benchmark refinement is non-decreasing (increasing only
  when the refined capability is maximal in $\dc G$); discovering subsumption
  structure can decrease $\volL$ by correcting overcounts
  (Proposition~\ref{prop:bootstrap_convergence}). This contraction is
  epistemically desirable: the actor's previous estimate was wrong, and
  structural discovery corrects it toward the true volume given full structural
  knowledge. The long-run trajectory converges toward this corrected value, but
  the convergence rate is not characterized.
\end{itemize}

\subsection{Open Problems Remaining}

\begin{itemize}
\item \textbf{Plurality:} Can GFM actors with different lattices (different
  capability descriptions, different benchmarks, different subsumption
  relations) coordinate? The lattice construction makes this more concrete:
  lattice reconciliation is a well-studied problem in knowledge representation
  \citep{ganter1999formal}.
\item \textbf{B-to-C gap:} The lattice measures capabilities (Candidate B)
  at binary-or-better resolution. The gap between capability volume and
  experiential optionality (Candidate C) is narrowed by the binary default
  (no known capability has zero weight) but not closed: the actor still
  understates capabilities whose latent gradations are unmeasured. Closing
  the gap requires developing richer benchmarks for experiential and
  relational capabilities.
\item \textbf{Continuous refinement:} As benchmark granularity increases ($\smax
  \to \infty$), does $\volL$ converge to a continuous measure? Under what
  conditions does the lattice discretization introduce errors that vanish in
  the limit?
\item \textbf{Lattice scaling:} The $O(n^3)$ computation is polynomial in
  $|\C|$, but $|\C|$ grows with population size, observation duration, and
  communication volume. A formal characterization of how $n$ scales with
  population parameters, and when sparsity reduces the effective cost, would
  strengthen the tractability claim.
\item \textbf{Automated benchmark refinement:} Can the GFM actor propose
  graded benchmarks for capabilities currently at the binary default ($w = 1$)?
  Richer benchmarks increase $w(d)$ and may enable deferred subsumption edges
  to pass the non-negativity check, integrating previously unstructured atoms
  into the lattice hierarchy.
\item \textbf{Subsumption correction:} Can mistakenly added edges be removed
  without enabling strategic lattice manipulation?
\item \textbf{Structural-avoidance incentive:} Can the objective be augmented
  to reward structural discovery independently of its effect on $\volL$? A
  regularization term penalizing structural uncertainty, or an information-gain
  bonus for subsumption discovery, would counteract the breadth-over-depth bias
  without abandoning $\volL$ as the primary objective.
\end{itemize}
